\documentclass[UTF8,a4paper,12pt]{ctexart}

\usepackage{amsmath,amssymb,bm}
\usepackage{booktabs}
\usepackage{geometry}
\usepackage{longtable}
\usepackage{array}
\usepackage{enumitem}
\usepackage{microtype}
\usepackage{hyperref}

\geometry{left=2.6cm,right=2.6cm,top=2.5cm,bottom=2.5cm}
\setlength{\parindent}{2em}
\setlength{\parskip}{0.35em}
\renewcommand{\arraystretch}{1.35}
\setlist{nosep}
\hypersetup{hidelinks}

\title{方法与案例研究部分}
\author{}
\date{}

\begin{document}

\maketitle

\section{3 面向多线换乘站疏散的队列感知路径组织模型}

本章建立面向多线换乘站整体疏散的队列感知路径组织模型。模型章的任务不是介绍某一搜索算法，而是把研究对象、客流来源、共享设施容量、到达时刻队列和评价指标统一到同一疏散表达中。AA 在本文中属于第 4 章的求解方法；它服务于本章模型，而不是替代本章模型。

大型五线换乘站应急疏散的关键不在于是否存在通向出口的可达路径，而在于多线路站台出站流、列车到达释放流、站厅滞留流和换乘流会沿同一组楼扶梯、换乘通道、闸机和出口前区共同加载。换乘客流在本文中被作为解释全站疏散过程的重要机制变量：其路径通常跨越多个线路空间，经过的服务链更长，也更容易与其他来源客流在同一共享设施处形成到达叠加。因此，本文不把研究目标缩小为换乘客流单独疏散，而是在整体疏散框架下刻画换乘流参与设施竞争后对关键设施队列、尾部清空和出口利用的影响。

\subsection{3.1 多线换乘站疏散服务链网络}

将车站公共区抽象为有向设施网络 \(G=(V,E)\)。节点集合 \(V\) 表示站台、站厅、换乘通道连接点、闸机区、出口前区和安全出口等空间单元；边集合 \(E\) 表示乘客可行走或可通过的连接关系。与一般最短路网络不同，本文把乘客从来源位置到安全出口的过程视为站内疏散服务链：不同来源客流在链上依次经过水平通道、垂向设施、服务设施和出口设施。对于楼梯、扶梯、通道、闸机和出口等具有明确服务能力的设施，定义共享资源集合 \(R\)。当多条边在网络表达中指向同一物理设施时，通过映射 \(\rho:E\rightarrow R\cup\varnothing\) 将其统一到同一资源容量下。

客流来源集合记为
\[
\mathcal G=\mathcal G^{out}\cup \mathcal G^{train}\cup \mathcal G^{hall}\cup \mathcal G^{tr},
\]
其中，\(\mathcal G^{out}\) 为站台等待或普通出站客流，\(\mathcal G^{train}\) 为列车到达释放客流，\(\mathcal G^{hall}\) 为站厅滞留客流，\(\mathcal G^{tr}\) 为换乘客流。对于换乘客流，模型保留其来源线路、目标方向或最终离站目标，使其通过量、等待时间和瓶颈暴露能够与其他来源组分开统计，但综合评价仍以全站乘客完成疏散为目标。

离散仿真时间集合为
\[
\mathcal T=\{0,\Delta t,2\Delta t,\ldots,T\}.
\]
在任一决策时刻 \(t\)，可参与路径组织的客流批次集合记为 \(\mathcal B(t)\)。一个批次 \(b\in\mathcal B(t)\) 表示为
\[
b=(g_b,o_b,D_b,n_b,t_b),
\]
其中，\(g_b\in\mathcal G\) 为来源组，\(o_b\in V\) 为批次当前所在节点或队列入口，\(D_b\subseteq V\) 为可接受安全出口集合，\(n_b\) 为批次人数，\(t_b\) 为批次进入路径组织阶段的时刻。批次建模的作用是保留列车到达和换乘释放的时间结构，也便于在路径选择后记录该批次对共享设施的未来到达需求。

对批次 \(b\)，设 \(\mathcal P_b(t)\) 为从 \(o_b\) 到 \(D_b\) 中任一出口的候选可行路径集合。路径选择变量定义为
\[
x_{bp}(t)=
\begin{cases}
1, & \text{时刻 }t\text{ 批次 }b\text{ 选择路径 }p,\\
0, & \text{否则}.
\end{cases}
\]
设施资源 \(r\in R\) 的服务率记为 \(\mu_r\)，时间步内可释放能力为 \(C_r(t)\)；资源入口等待队列为 \(Q_r(t)\)。空间节点 \(v\in V\) 的瞬时占用人数为 \(N_v(t)\)，最大接收能力为 \(\bar N_v\)。边 \(e\in E\) 上处于行走或过渡状态的人数记为 \(M_e(t)\)。已完成疏散人数为 \(E^{out}(t)\)，批次 \(b\) 的实际离站时刻记为 \(T_b\)。

\begin{table}[htbp]
\centering
\caption{路径组织模型的主要符号}
\small
\begin{tabular}{p{0.18\textwidth}p{0.70\textwidth}}
\toprule
符号 & 含义 \\
\midrule
\(V,E,R\) & 站内节点、可通行边和共享设施资源集合 \\
\(\rho(e)\) & 边 \(e\) 对应的物理共享设施资源；若不对应资源则为空 \\
\(\mathcal G\) & 客流来源集合，包括站台出站、列车到达、站厅滞留和换乘客流 \\
\(\mathcal B(t)\) & 时刻 \(t\) 进入路径组织阶段的客流批次集合 \\
\(\mathcal P_b(t)\) & 批次 \(b\) 在时刻 \(t\) 的候选可行路径集合 \\
\(x_{bp}(t)\) & 批次 \(b\) 是否选择路径 \(p\) 的决策变量 \\
\(Q_r(t),\mu_r,C_r(t)\) & 共享设施 \(r\) 的等待队列、服务率和时间步释放能力 \\
\(N_v(t),\bar N_v\) & 节点 \(v\) 的占用人数和最大接收能力 \\
\(T_b,E^{out}(t)\) & 批次 \(b\) 的实际离站时刻和时刻 \(t\) 的累计疏散人数 \\
\bottomrule
\end{tabular}
\end{table}

\subsection{3.2 共享容量约束下的路径组织优化模型}

多线换乘站疏散路径组织的核心矛盾在于：乘客在几何上可能存在多条可达出口路径，但这些路径会在若干楼扶梯、通道、闸机和出口前区处共享同一物理服务能力。若路径选择只根据当前边权或当前节点状态进行评价，就难以反映已经被前序批次承诺占用的未来容量，也难以解释为什么几何上较短的路径可能在乘客真正到达设施时转化为长时间等待。为此，本文将路径组织写成带共享容量约束的动态离散优化问题。

模型目标包括五类：缩短整体疏散时间，降低共享设施累计排队，减少空间阻塞和停滞，避免无必要绕行，并控制拥挤暴露。设 \(L_p\) 为路径 \(p\) 的几何长度，\(S_v(t)\) 为节点 \(v\) 的空间阻塞或停滞量，\(R_{bp}(t)\) 为批次 \(b\) 选择路径 \(p\) 时的拥挤暴露辅助项，则目标函数表示为
\[
\min Z=Z_T+\alpha Z_Q+\beta Z_S+\gamma Z_L+\lambda Z_R.
\]
其中，
\[
Z_T=\sum_{b} n_b T_b,
\]
\[
Z_Q=\sum_{t\in\mathcal T}\sum_{r\in R} Q_r(t)\Delta t,
\]
\[
Z_S=\sum_{t\in\mathcal T}\sum_{v\in V} S_v(t)\Delta t,
\]
\[
Z_L=\sum_{t\in\mathcal T}\sum_{b\in\mathcal B(t)}\sum_{p\in\mathcal P_b(t)}n_b L_p x_{bp}(t),
\]
\[
Z_R=\sum_{t\in\mathcal T}\sum_{b\in\mathcal B(t)}\sum_{p\in\mathcal P_b(t)}n_b R_{bp}(t)x_{bp}(t).
\]
权重 \(\alpha,\beta,\gamma,\lambda\) 用于把等待、阻塞、绕行和拥挤暴露统一到路径组织目标中。本文不把这些权重解释为乘客心理偏好参数，而将其作为路径组织模型中不同代价项的尺度控制参数；具体取值应与当前模型配置和后续敏感性检查对应。

每个可决策批次在一个决策时刻只能选择一条执行路径：
\[
\sum_{p\in\mathcal P_b(t)}x_{bp}(t)=1,\quad b\in\mathcal B(t).
\]
若路径不连通、出口不可用、通过受限设施不可达，或下游空间在预测期内不满足接收条件，则该路径不进入 \(\mathcal P_b(t)\)，等价地有
\[
x_{bp}(t)=0,\quad p\notin\mathcal P_b(t).
\]
设 \(I_{brp}(t,\theta)\) 表示批次 \(b\) 在时刻 \(t\) 选择路径 \(p\) 后是否预计于时间区间 \([\theta,\theta+\Delta t)\) 到达资源 \(r\)。则资源 \(r\) 在时间区间内承受的预测到达需求为
\[
A_r(\theta)=\sum_{t\leq \theta}\sum_{b\in\mathcal B(t)}\sum_{p\in\mathcal P_b(t)} n_b x_{bp}(t)I_{brp}(t,\theta).
\]
资源时间步服务能力由服务率和离散时间步长共同决定：
\[
C_r(t)=\left\lfloor \mu_r\Delta t+\xi_r(t)\right\rfloor,
\]
其中，\(\xi_r(t)\) 为容量离散化中的余量项。资源入口队列按照到达需求和服务释放更新：
\[
Q_r(t+\Delta t)=\max\{Q_r(t)+A_r(t)-C_r(t),0\}.
\]
这一约束的含义不是形式上让路径层和执行层“保持一致”，而是保证同一楼梯、扶梯、闸机或出口在一个时间步内只释放一次真实容量。多条路径同时使用同一设施时，所有来源客流必须在同一容量约束下竞争，而不能因为网络中存在多条边而重复获得服务能力。

除设施服务能力外，路径组织还受下游空间接收能力约束。设 \(B_v(t)\) 为当前时间步预计进入节点 \(v\) 的人数，则
\[
N_v(t)+B_v(t)\leq \bar N_v,\quad v\in V.
\]
当该约束不能满足时，乘客不能无条件进入下游空间，而应在上游节点、边末端或设施入口形成等待。该约束用于刻画站厅、通道、闸机前区和出口前区在高负荷下可能出现的空间接收受限过程。

系统还需满足人数守恒：
\[
\sum_{v\in V}N_v(t)+\sum_{e\in E}M_e(t)+\sum_{r\in R}Q_r(t)+E^{out}(t)=N^{tot},
\]
其中，\(N^{tot}\) 为场景初始总人数。已经进入边上移动过程、已经排队等待服务、已经停留在节点中和已经离站的乘客共同构成全系统人数，仿真过程中不应出现人数丢失或重复计数。

最后，路径组织具有非预知性。时刻 \(t\) 的决策只能使用当前可获得状态 \(\Omega(t)\)，包括节点占用、资源队列、在途客流、已登记未来到达事件和可用出口状态；不能使用执行层未来实际尚未产生的结果。已进入执行过程的客流批次，其在途状态和已承诺资源到达事件不得被后续批次任意改写。该约束保证模型反映真实应急疏散中的滚动决策过程，而不是事后用全局信息重新安排所有乘客路径。

\subsection{3.3 共享设施到达时刻队列预测模型}

共享设施队列预测模型为路径组织模型提供关键输入。它回答的问题不是“当前设施是否拥堵”，而是“某一批乘客真正到达该设施时，前序批次和在途客流是否已经对该设施形成排队压力”。这一区别是本文模型与只使用当前状态边权的路径组织方法之间的主要差异。

设当前时刻为 \(t\)，资源 \(r\) 当前入口队列为 \(Q_r(t)\)。在路径组织过程中，系统保存两类未来到达事件：一类来自已经进入执行过程、尚未完成资源服务或尚在通行过程中的实际在途客流；另一类来自当前决策轮次内前序批次已经选择路径但尚未执行的路径意图。资源 \(r\) 的未来到达事件集合记为
\[
\mathcal A_r(t)=\{(\eta_j,a_j,g_j,b_j)\mid \eta_j\geq t\},
\]
其中，\(\eta_j\) 为事件到达资源入口的预计时刻，\(a_j\) 为到达人数，\(g_j\) 为来源组，\(b_j\) 为批次标识。该事件集合使同一时间步内先被组织的客流批次能够影响后续批次对同一资源未来等待时间的判断。

对任一候选到达时刻 \(\theta\geq t\)，预测队列可写为
\[
\widehat Q_r(\theta|t)=
\max\left\{
Q_r(t)+\sum_{(\eta_j,a_j)\in\mathcal A_r(t),\eta_j\leq \theta}a_j
-\mu_r(\theta-t),0
\right\}.
\]
据此得到批次到达时的预测等待时间：
\[
\widehat W_r(\theta|t)=\frac{\widehat Q_r(\theta|t)}{\mu_r}.
\]
对于批次自身进入资源后占用服务能力造成的平均服务时间，采用
\[
S_r(n_b)=\frac{n_b}{2\mu_r}.
\]
该项用于近似同一批次内部乘客在设施前后相继通过时形成的平均服务占用；它不是额外人为惩罚，而是把批次规模转化为设施服务过程中的时间代价。

给定候选路径 \(p=(e_1,e_2,\ldots,e_m)\)，从当前时刻 \(t\) 开始递推批次沿路径的预计到达时刻。设进入边 \(e_i\) 或其对应资源的时刻为 \(\eta_i\)，边 \(e_i\) 的物理行走时间为 \(\tau_{e_i}(\eta_i)\)，则路径预测广义代价为
\[
\widehat C_{bp}(t)=
\sum_{i=1}^{m}
\left[
\tau_{e_i}(\eta_i)+
I_{\rho(e_i)\in R}
\left(
\widehat W_{\rho(e_i)}(\eta_i|t)+S_{\rho(e_i)}(n_b)
\right)
\right]+
\lambda R_{bp}(t).
\]
其中，\(I_{\rho(e_i)\in R}\) 用于判断边 \(e_i\) 是否对应共享设施资源。若一条路径几何距离较短，但其关键共享设施在该批次到达时已经被前序客流占用，则该路径的预测代价可能高于几何距离更长但未来队列更短的路径。这样，模型把路径长度、到达时刻等待和设施服务过程放入同一时间维度评价。

需要强调的是，到达时刻队列预测并不等同于提前知道未来真实状态。它只使用当前已经观测到的状态、在途客流和本轮已经生成的路径意图进行推算。后续执行层仍按照真实共享容量释放乘客；若实际排队和预测不一致，下一时间步会根据新的 \(\Omega(t+\Delta t)\) 重新组织路径。

\subsection{3.4 车站疏散效果综合评价方法}

为了避免把论文写成单纯的算法快慢比较，本文评价体系围绕全站疏散效率、共享设施瓶颈、来源组差异和跨模型趋势一致性展开。整体效率指标包括 \(T_{50}\)、\(T_{80}\)、\(T_{95}\)、\(T_{100}\)、平均疏散时间和累计停滞人秒。其中，
\[
T_{\theta}=\min\{t\mid E^{out}(t)\geq \theta N^{tot}\},\quad \theta\in\{0.50,0.80,0.95,1.00\}.
\]
\(T_{95}\) 和 \(T_{100}\) 用于刻画大部分乘客离站后的尾部清空阶段，避免仅用平均疏散时间掩盖少数滞留乘客对安全组织的影响。

共享设施瓶颈指标包括资源最大队列、累计队列人秒、资源利用率和队列持续时间。累计队列人秒记为
\[
\Phi_Q=\sum_{t\in\mathcal T}\sum_{r\in R}Q_r(t)\Delta t.
\]
若某一资源的 \(\Phi_Q\)、峰值队列和持续时间同时较高，说明其不是短时波动，而是对疏散过程产生连续约束的关键设施。空间阻塞指标记为
\[
\Phi_S=\sum_{t\in\mathcal T}\sum_{v\in V}S_v(t)\Delta t,
\]
用于识别下游空间接收不足造成的停滞和排队扩散。

来源组评价重点关注换乘客流如何影响整体疏散，而不是将换乘客流从全站目标中分离出去。对来源组 \(g\)，统计其离站曲线 \(E_g(t)\)、清空时间 \(T_g\)、平均等待时间和关键设施暴露。若换乘客流在高负荷场景中更集中地占用某些楼扶梯、通道或闸机资源，则应进一步分析该资源是否同时服务其他线路出站客流，以及这种叠加是否改变全站尾部清空。

Pathfinder 微观仿真对照用于检查宏观或中观图模型结论在微观行人仿真中的趋势表现。由于两类模型的行人行为假设、空间表达和冲突处理机制不同，本文不把 Pathfinder 结果写成逐人严格校准，而将其作为趋势对照，重点比较疏散曲线形态、拥堵位置、出口利用分布和关键来源组滞留特征。

\section{4 队列感知路径组织模型的求解方法}

第 3 章给出了面向多线换乘站疏散的队列感知路径组织模型，包括服务链网络、路径组织目标、共享容量约束、到达时刻队列预测和评价指标。本章说明该模型如何转化为可在离散疏散仿真中执行的求解过程。完整模型具有动态、离散、共享容量和非预知性等特点，若在每一时刻对所有批次、所有路径和所有未来状态做全局整数优化，计算规模会迅速扩大，也难以服务应急疏散中的滚动更新。因此，本文采用滚动动态加载转换，将全局路径组织问题分解为连续时间步内的批次路径搜索和共享容量执行。在求解层，A* 提供基本的启发式图搜索框架，AA 则在该框架上引入到达时刻队列预测和多标签状态表达，使路径搜索能够处理共享设施未来排队对路径代价的影响。

\subsection{4.1 共享容量约束下的滚动动态加载转换}

在每个离散时间步开始时，系统读取上一时间步结束后的状态并冻结为当前决策状态：
\[
\Omega(t)=\{N_v(t),Q_r(t),M_e(t),E^{out}(t),\mathcal A_r(t)\}.
\]
状态冻结的作用是使同一时间步内所有批次从同一执行状态出发进行路径组织，避免由于批次处理顺序不同而随意改变节点占用、资源队列或下游接收能力。随后，系统识别当前可调度批次集合 \(\mathcal B(t)\)，包括已经位于节点等待区的乘客、到达设施入口但尚未获得服务的乘客，以及在当前时刻释放到站内网络中的列车到达客流。

对每个批次 \(b\in\mathcal B(t)\)，求解器在候选路径集合 \(\mathcal P_b(t)\) 上求解滚动子问题：
\[
p_b^{*}(t)=\arg\min_{p\in\mathcal P_b(t)}
\widehat C_{bp}(t|\Omega(t),\mathcal A(t)).
\]
其中，\(\widehat C_{bp}\) 由第 3.3 节的到达时刻队列预测模型给出。路径 \(p_b^{*}(t)\) 被接受后，系统不是立即把所有乘客强行移动到出口，而是把该路径上预计到达共享资源的时间和人数写入本轮未来到达事件集合：
\[
\mathcal A_r(t)\leftarrow
\mathcal A_r(t)\cup\{(\widehat\eta_{brp},n_b,g_b,b)\}.
\]
这样，后续批次在同一轮路径组织中可以看到前序批次对未来资源容量的占用。该机制对应的是路径组织层的承诺到达，而不是执行层已经完成的实际通行；实际通行仍由共享容量执行层决定。

当本时间步路径意图生成完毕后，执行层按照资源服务能力和下游接收能力释放乘客。若资源 \(r\) 在当前时间步可服务 \(C_r(t)\) 人，则所有等待进入该资源的来源客流共同竞争这一容量。若下游节点达到接收上限，则即使资源有剩余服务能力，乘客也不能继续向下游无限堆积，而应在上游形成等待。执行完成后更新 \(N_v(t+\Delta t)\)、\(Q_r(t+\Delta t)\)、\(M_e(t+\Delta t)\) 和 \(E^{out}(t+\Delta t)\)，进入下一时间步重新求解。

上述转换的意义在于把第 3 章的容量约束、队列约束和非预知性约束落到可执行流程中。路径组织层只生成意图，执行层负责验证这些意图在真实共享容量下能否释放。二者分离但共用同一网络、同一设施容量和同一客流状态，因此可以避免同一物理设施被不同路径重复计入容量。

\subsection{4.2 A* 基础搜索框架}

A* 是在加权图上求解最小代价路径的启发式搜索方法。本文介绍 A* 的目的不是展开通用算法教材，而是说明 AA 在哪几个环节上对基础 A* 进行扩展。对处于节点 \(v\) 的搜索状态，A* 用评价函数
\[
f(v)=g(v)+h(v)
\]
确定优先扩展顺序。其中，\(g(v)\) 表示从起点到节点 \(v\) 已经累计的实际代价，\(h(v)\) 表示从节点 \(v\) 到目标节点的剩余代价估计。若 \(h(v)\) 是不超过真实剩余代价的下界，A* 可以在减少无效扩展的同时保持最优路径搜索的合理性。

在静态路径规划中，边代价通常由距离、速度或固定通行阻抗给出。此时，同一个节点只需要保留当前最小的 \(g(v)\)，因为到达该节点后面对的后续边代价不会因到达时刻不同而变化。地铁站疏散路径组织与此不同：乘客到达楼扶梯、通道、闸机或出口前区的时刻会影响其面对的队列长度和等待时间。也就是说，同一节点在不同时刻被到达时，后续路径代价可能不同。若仍按静态 A* 将节点状态简单合并，可能删除一条当前代价略高但后续等待更低的可行路径。

因此，本文没有把 A* 作为完整疏散模型，而是将其作为滚动子问题的搜索骨架。普通 A* 中的优先队列、累计代价、启发式下界和前驱回溯机制仍然保留；需要改造的是状态表达和边代价计算。对批次 \(b\)，本文采用自由流条件下到可用出口集合 \(D_b\) 的最短剩余时间作为启发式下界：
\[
h(v)=\frac{d^{free}(v,D_b)}{v^{free}}.
\]
该下界只反映在无排队、无阻塞条件下从当前节点到出口的最短可能时间，不包含未来队列等待。因此，它不会替代实际路径代价，而是用于限制搜索空间。实际边代价则在 AA 中根据乘客预计进入该边或共享设施的时刻动态计算。

\subsection{4.3 AA 时间依赖多标签路径搜索}

本文将用于求解滚动子问题的时间依赖路径搜索方法记为 AA，即 AdaptiveQueueAwareAStar。该方法的作用不是重新定义第 3 章的模型，而是在 A* 搜索框架中调用到达时刻队列预测，使搜索标签能够感知不同到达时间对应的设施等待差异。与只改变边权的改进 A* 不同，AA 的关键在于把节点状态从单一空间位置扩展为空间位置、到达时刻、累计代价和拥挤暴露共同决定的搜索标签。

AA 对批次 \(b\) 从当前节点 \(o_b\) 到目标出口集合 \(D_b\) 进行多标签 A* 搜索。搜索标签记为
\[
\ell=(v,\eta,c,\psi,\pi),
\]
其中，\(v\) 为标签所在节点，\(\eta\) 为预计到达该节点的时刻，\(c\) 为从起点到当前节点的累计广义代价，\(\psi\) 为累计拥挤暴露或风险代价，\(\pi\) 为前驱标签。与静态最短路不同，同一节点上的两个标签即使位置相同，只要到达时刻不同，也可能面对不同的未来设施队列，因此不能简单按节点合并。AA 保留必要的时间标签，并通过支配关系删除明显劣势标签：若两个标签位于同一节点，且一个标签到达不晚、累计代价和拥挤暴露均不高，则另一个标签可被剔除。

当标签沿边 \(e=(u,v)\) 扩展时，首先计算物理行走时间 \(\tau_e(\eta)\)。若边 \(e\) 不对应共享资源，则新标签到达时刻和代价主要由行走时间更新；若 \(\rho(e)=r\)，则在预计进入资源的时刻 \(\eta\) 查询 \(\widehat Q_r(\eta|t)\)，并加入预测等待 \(\widehat W_r(\eta|t)\) 和批次自身服务占用 \(S_r(n_b)\)。若下游节点在预测期内存在接收受限，则扩展代价还应包含相应的空间等待或阻塞代价。标签扩展可概括为
\[
c'=
c+\tau_e(\eta)+I_{\rho(e)\in R}
\left(\widehat W_{\rho(e)}(\eta|t)+S_{\rho(e)}(n_b)\right)
\;+
\lambda R_e(\eta).
\]
搜索使用自由流条件下到最近可用出口的时间下界作为启发式函数：
\[
h(v)=\frac{d^{free}(v,D_b)}{v^{free}}.
\]
该启发式只用于减少搜索空间，不替代真实路径代价。真实路径代价仍由已递推的行走时间、预测设施等待、批次服务占用和拥挤暴露共同确定。

AA 输出的路径不仅包括节点序列，还包括该批次预计到达各共享设施的时间、人数和来源组信息。若该路径被滚动子问题接受，这些信息会登记到当前时间步的未来到达事件集合，使随后处理的批次能够感知同一设施即将承受的负荷。由此，AA 的“自适应”不是泛泛地根据拥堵调权，而是随前序批次的路径承诺更新到达时刻队列预测，再据此重新评价后续批次路径。该处理与当前实现保持一致：搜索默认采用多标签模式，启发式为自由流剩余时间下界，资源等待由预测队列和服务率计算，路径接受后只登记共享设施的软规划到达事件，实际通行仍交给共享容量执行层完成。

\subsection{4.4 共享容量执行与计算对照方法}

AA 搜索得到的是路径组织意图，不能直接等同于乘客已经按该路径完成疏散。因此，所有路径组织结果都必须进入统一的共享容量执行层。执行层按照第 3.2 节的资源容量约束和下游接收约束推进乘客，记录实际队列、实际通行、空间阻塞、出口通过量和来源组清空时间。共享容量执行层的意义在于容量守恒和同步竞争：同一物理楼梯、扶梯、闸机或出口前区在同一时间步只释放一次服务能力，所有来源客流按照同一约束竞争该能力。若路径组织阶段对某设施的未来负荷判断准确，执行层中该设施的实际队列应与预测趋势相符；若出现偏差，下一时间步会根据更新后的状态重新组织。

本文使用 ImprovedAStar 作为计算对照方法。两种方法使用相同的站内网络、相同的客流输入、相同的速度模型、相同的共享设施容量、相同的下游接收约束和相同的终止条件。二者差异限定在路径组织层：AA 在搜索过程中登记并查询未来到达事件，从而评价批次到达设施时的预测队列；ImprovedAStar 使用当前网络状态和即时边权进行路径评价，不引入本轮已经承诺的未来到达负荷。ImprovedAStar 仅是计算基准，不对应任何已核验的现场疏散安排，也不承担现实预案的角色。

求解流程可概括为以下步骤：
\begin{enumerate}
\item 读取并冻结当前时间步的节点占用、设施队列、在途客流和已登记未来到达事件；
\item 识别当前可参与路径组织的客流批次，并生成其候选出口路径集合；
\item 对每个批次调用 AA 时间依赖 A* 搜索，计算包含行走时间、到达时刻等待、批次服务占用和拥挤暴露的路径代价；
\item 接受路径后登记该批次对共享设施的预计到达事件，使后续批次能够感知前序批次对未来容量的占用；
\item 将全部路径意图交给共享容量执行层，在资源服务率和下游接收能力约束下释放乘客；
\item 统计队列、停滞、出口利用和来源组清空状态，进入下一时间步滚动求解，直至所有乘客完成疏散。
\end{enumerate}

上述流程将模型目标、预测模块和执行约束连接起来。若只报告 AA 与 ImprovedAStar 的总疏散时间差异，无法说明差异来自何处；因此案例研究需要先分析到达负荷与队列预测结果，再分析路径组织方案性能，最后将其放回完整疏散执行和 Pathfinder 微观仿真中检查。

\section{5 龙阳路站案例研究}

本文以龙阳路站作为案例，检验所建模型在大型多线换乘站整体应急疏散中的适用性。案例站连接多条城市轨道交通线路和磁浮线路，站内包括多个线路站台、站厅、换乘通道、楼扶梯、闸机区和地面出口。该类车站的关键特征不是单一路径是否可达，而是多线路客流、列车到达客流和换乘客流会在若干共享设施上叠加，使设施到达时刻队列成为影响全站尾部清空的重要因素。

案例建模使用当前已经能够追溯的研究材料：CAD 图纸和节点边关系用于构建站内疏散网络，当前模型配置用于生成低负荷和高负荷客流输入，AA 与 ImprovedAStar 用于同一执行层下的路径组织对照，Pathfinder 用于微观趋势对照。设施服务能力和行为参数在定稿时应在参数表中区分 CAD 可量测值、模型设定值、文献取值和必要假设值；尚未完成来源核验的值不写成实测值。速度-密度关系仅作为基础运动时间估计的一部分，本文实验不把它作为主要敏感性分析对象。

案例设置低负荷常规应急场景和高负荷满载列车叠加场景。根据当前模型配置，低负荷场景包含站台等待、站厅滞留和换乘客流，共 2187 人；高负荷场景在此基础上叠加多线路列车到达释放客流，共 17905 人。上述人数为当前模型配置值，若后续仿真输入冻结时发生修改，应在正式稿中同步更新。

案例研究的验证对象不是某一算法名称本身，而是第 3 章提出的“到达时刻队列感知路径组织”机制。尚未完全冻结的实验输出只保留表格结构和解释口径，不提前写入定量结论。

\begin{table}[htbp]
\centering
\caption{案例场景设置}
\small
\begin{tabular}{p{0.23\textwidth}p{0.32\textwidth}p{0.33\textwidth}}
\toprule
场景 & 客流组成 & 实验目的 \\
\midrule
低负荷常规应急场景 & 站台等待客流、站厅滞留客流和换乘客流；当前模型配置总人数为 2187 人 & 检查在共享设施竞争不强时，AA 是否保持稳定路径组织，避免因预测机制产生无必要绕行或额外等待 \\
高负荷满载列车叠加场景 & 低负荷基础客流叠加多线路列车到达释放客流；当前模型配置总人数为 17905 人 & 放大多源客流到达重叠和共享设施容量竞争，检验到达时刻队列预测对尾部清空和关键设施负荷分布的作用 \\
\bottomrule
\end{tabular}
\end{table}

案例研究按三步展开。第一步分析到达时刻负荷与队列预测结果，用于说明 AA 的中间预测机制是否确实识别了关键共享设施的未来排队压力。第二步比较 AA 与 ImprovedAStar 的路径组织方案求解性能，重点考察路径代价、设施负荷分布、尾部清空和计算耗时。第三步将路径组织结果放回完整疏散执行过程，并结合 Pathfinder 微观仿真和敏感性分析，讨论结论在不同换乘客流比例、列车到达重叠程度和关键设施服务能力下的适用边界。

\subsection{5.1 到达时刻负荷与队列预测结果分析}

案例研究首先分析到达负荷与队列预测结果，而不是直接进入总疏散时间比较。这样安排的原因是：本文模型的核心并非单纯寻找几何最短路径，而是预测乘客批次到达共享设施时的队列等待。如果这一中间机制不能被说明，那么后续 AA 与 ImprovedAStar 的差异就容易被误读为一般性的算法性能差异。

本节选取若干代表性共享设施，包括高换乘参与度的楼扶梯、连接不同线路空间的换乘通道、出站闸机组和出口前区。对每一类设施，比较三类时间序列：当前状态下可见队列、AA 预测的到达时刻队列，以及执行层形成的实际队列。低负荷场景中，重点观察预测队列是否接近零或保持短时波动；高负荷场景中，重点观察预测模型是否能在实际队列峰值出现前识别到达高峰。

\begin{table}[htbp]
\centering
\caption{到达时刻队列预测结果记录表}
\small
\begin{tabular}{p{0.18\textwidth}p{0.16\textwidth}p{0.17\textwidth}p{0.18\textwidth}p{0.21\textwidth}}
\toprule
关键共享设施 & 场景 & 预测到达峰值时段 & 执行层队列峰值 & 解释重点 \\
\midrule
楼扶梯资源 A & 低负荷/高负荷 & 待冻结 & 待冻结 & 判断换乘流与站台流是否在该设施形成到达叠加 \\
换乘通道资源 B & 低负荷/高负荷 & 待冻结 & 待冻结 & 判断跨线服务链是否成为尾部清空的中间约束 \\
闸机组资源 C & 低负荷/高负荷 & 待冻结 & 待冻结 & 判断多来源客流离站前是否在出站设施处汇聚 \\
出口前区资源 D & 低负荷/高负荷 & 待冻结 & 待冻结 & 判断出口利用是否受上游队列转移影响 \\
\bottomrule
\end{tabular}
\end{table}

若结果显示高负荷场景中 AA 在路径选择阶段已经预测到部分共享设施的到达峰值，而 ImprovedAStar 在同一时刻只反映当前队列，则可说明二者差异来自对未来设施负荷的处理方式，而不是简单的参数调权。若低负荷场景中预测队列和实际队列均较低，则说明该场景主要用于检查 AA 的稳定性：在设施竞争不明显时，模型不应为了追求表面分流而显著增加路径长度或等待时间。

\subsection{5.2 路径组织方案求解性能分析}

在到达负荷与队列预测结果得到说明后，本节比较 AA 与 ImprovedAStar 的路径组织方案性能。比较采用相同的站内网络、客流输入和共享容量执行层，因此结果差异应主要来自路径组织层对未来设施队列的处理方式。低负荷场景侧重观察 AA 是否与对照方法保持接近的总体效率，并避免不必要绕行；高负荷场景侧重观察 AA 是否改变关键共享设施的负荷分布，从而影响停滞等待和尾部清空。

\begin{table}[htbp]
\centering
\caption{路径组织方案总体性能比较}
\small
\begin{tabular}{p{0.16\textwidth}p{0.15\textwidth}p{0.13\textwidth}p{0.13\textwidth}p{0.13\textwidth}p{0.13\textwidth}p{0.15\textwidth}}
\toprule
场景 & 方法 & \(T_{50}\) & \(T_{80}\) & \(T_{95}\) & \(T_{100}\) & 解释重点 \\
\midrule
低负荷常规应急场景 & ImprovedAStar & 待冻结 & 待冻结 & 待冻结 & 待冻结 & 计算对照基准 \\
低负荷常规应急场景 & AA & 待冻结 & 待冻结 & 待冻结 & 待冻结 & 检查稳定性与非瓶颈条件下的额外代价 \\
高负荷满载列车叠加场景 & ImprovedAStar & 待冻结 & 待冻结 & 待冻结 & 待冻结 & 识别当前状态路径组织下的尾部清空表现 \\
高负荷满载列车叠加场景 & AA & 待冻结 & 待冻结 & 待冻结 & 待冻结 & 检查未来队列预测是否改善尾部清空和等待暴露 \\
\bottomrule
\end{tabular}
\end{table}

仅比较 \(T_{100}\) 不足以解释路径组织机制。因此，本节还应分解累计等待和空间阻塞。对每个场景和方法，统计资源队列人秒、空间阻塞人秒、累计停滞时间、总移动距离和计算耗时。若 AA 在高负荷场景中降低某些关键资源的队列峰值，却显著增加总移动距离，则需要进一步判断这种绕行是否换来了尾部清空改善；若总疏散时间改善不明显，但关键设施峰值降低，也应从安全组织和拥堵扩散角度解释其意义。

\begin{table}[htbp]
\centering
\caption{路径组织代价分解与瓶颈指标}
\small
\begin{tabular}{p{0.14\textwidth}p{0.13\textwidth}p{0.14\textwidth}p{0.14\textwidth}p{0.14\textwidth}p{0.13\textwidth}p{0.13\textwidth}}
\toprule
场景 & 方法 & 队列人秒 & 空间阻塞人秒 & 累计停滞时间 & 总移动距离 & 计算耗时 \\
\midrule
低负荷 & ImprovedAStar & 待冻结 & 待冻结 & 待冻结 & 待冻结 & 待冻结 \\
低负荷 & AA & 待冻结 & 待冻结 & 待冻结 & 待冻结 & 待冻结 \\
高负荷 & ImprovedAStar & 待冻结 & 待冻结 & 待冻结 & 待冻结 & 待冻结 \\
高负荷 & AA & 待冻结 & 待冻结 & 待冻结 & 待冻结 & 待冻结 \\
\bottomrule
\end{tabular}
\end{table}

此外，应输出关键设施负荷分布。对于每个代表性资源，统计不同来源组的到达人数、通过人数、峰值队列和队列持续时间。该分析用于回答换乘客流在整体疏散中的机制作用：若某些共享设施的队列主要由换乘流与列车到达流叠加形成，则 AA 的路径调整应解释为对多服务链交叉加载的响应，而不是对某一类乘客单独优先。

\subsection{5.3 路径组织与全站疏散效果集成分析}

本节将路径组织方案放回完整疏散执行过程，分析路径组织、共享设施队列和乘客疏散效果之间的关系。与只展示最终疏散时间不同，集成分析需要说明三件事：第一，AA 是否改变了多来源客流在站内服务链上的加载顺序和设施选择；第二，这种变化是否在执行层表现为关键设施排队、停滞等待或出口利用的改变；第三，Pathfinder 微观仿真是否在拥堵位置和疏散曲线趋势上支持图模型结论。

\subsubsection{5.3.1 多来源客流与换乘流机制分析}

对站台出站、列车到达、站厅滞留和换乘客流分别统计离站曲线和关键设施暴露。换乘客流不被写成本文唯一优化对象，而作为解释共享设施竞争的重要来源组。若高负荷场景中换乘客流跨越较长服务链，并与列车到达释放客流在相同楼扶梯、通道或闸机处形成叠加，则需要说明 AA 是否通过改变部分批次路径，削弱某一共享设施的到达峰值，或将客流转移到尚有服务余量的出口方向。

\begin{table}[htbp]
\centering
\caption{来源组疏散效果与关键设施暴露}
\small
\begin{tabular}{p{0.17\textwidth}p{0.15\textwidth}p{0.15\textwidth}p{0.17\textwidth}p{0.25\textwidth}}
\toprule
来源组 & 方法 & 清空时间 & 平均等待时间 & 关键解释 \\
\midrule
站台出站客流 & ImprovedAStar/AA & 待冻结 & 待冻结 & 判断普通出站流是否因换乘流叠加而受阻 \\
列车到达释放客流 & ImprovedAStar/AA & 待冻结 & 待冻结 & 判断列车集中释放是否放大共享设施到达峰值 \\
站厅滞留客流 & ImprovedAStar/AA & 待冻结 & 待冻结 & 判断站厅空间接收约束是否影响后续离站 \\
换乘客流 & ImprovedAStar/AA & 待冻结 & 待冻结 & 判断跨线服务链是否参与关键瓶颈和尾部清空 \\
\bottomrule
\end{tabular}
\end{table}

\subsubsection{5.3.2 Pathfinder 微观仿真趋势对照}

Pathfinder 对照用于检验图模型输出在微观空间中的趋势表现。对照内容包括疏散曲线、拥堵位置、出口利用和来源组滞留特征。由于 Pathfinder 的行人避让、局部冲突和空间占用机制与图模型不同，本文不应写成逐个乘客轨迹完全一致，而应说明两类模型是否在关键瓶颈位置、拥堵持续时段和总体疏散曲线形态上保持一致。

\begin{table}[htbp]
\centering
\caption{图模型与 Pathfinder 对照维度}
\small
\begin{tabular}{p{0.20\textwidth}p{0.33\textwidth}p{0.35\textwidth}}
\toprule
对照维度 & 图模型输出 & Pathfinder 输出 \\
\midrule
疏散曲线 & 累计疏散人数、\(T_{95}\)、\(T_{100}\) & 个体离站时间分布和累计疏散曲线 \\
拥堵位置 & 关键设施队列峰值、队列人秒和空间阻塞人秒 & 密度云图、拥堵持续时段和局部聚集区域 \\
出口利用 & 各出口通过人数及来源组构成 & 各出口实际离站人数和占比 \\
来源组差异 & 换乘、站台、列车到达和站厅客流清空时间 & 不同人群组离站时间和拥堵暴露 \\
\bottomrule
\end{tabular}
\end{table}

若 Pathfinder 中的高密度区域与图模型识别的高队列资源一致，且疏散曲线尾部变化方向相同，则可认为微观仿真在趋势上支持路径组织模型结论。若两者不一致，应优先检查空间几何表达、局部避让机制、设施容量设定和出口边界条件，而不能直接把差异解释为某一方法失效。

\subsubsection{5.3.3 敏感性分析与适用边界}

敏感性分析应服务于本文模型机制，而不是泛泛改变所有行人参数。本文重点分析三类因素：换乘客流比例、列车到达重叠程度和关键共享设施服务能力。换乘客流比例用于检验服务链交叉强度变化时 AA 的作用边界；列车到达重叠程度用于检验多来源客流在时间上集中释放时到达时刻预测的价值；关键共享设施服务能力用于检验瓶颈强弱变化对路径组织效果的影响。

\begin{table}[htbp]
\centering
\caption{敏感性分析设置与解释目的}
\small
\begin{tabular}{p{0.24\textwidth}p{0.30\textwidth}p{0.34\textwidth}}
\toprule
敏感性因素 & 调整方式 & 解释目的 \\
\midrule
换乘客流比例 & 在总需求或基础需求不变条件下调整换乘来源组占比 & 检验服务链交叉增强时，AA 是否更能识别共享设施未来排队 \\
列车到达重叠程度 & 调整多线路列车释放的时间间隔或同步程度 & 检验客流到达高峰越集中时，到达时刻队列预测是否更重要 \\
关键共享设施服务能力 & 调整代表性楼扶梯、通道、闸机或出口前区容量 & 检验设施瓶颈强弱变化下，AA 的适用范围和失效边界 \\
\bottomrule
\end{tabular}
\end{table}

结果解释应围绕适用边界展开。若 AA 的优势主要出现在高换乘比例、高到达重叠或关键设施服务能力较低的条件下，应将其解释为面向共享设施竞争明显场景的路径组织方法；若在低负荷或瓶颈较弱场景中差异较小，则说明模型不会在非瓶颈条件下强行产生显著改善。这样的解释比单纯报告“AA 更快”更能说明方法适用于哪些多线换乘站疏散情境。

\end{document}
